\documentclass[a4paper,11pt]{report}
\usepackage[table]{xcolor}
\usepackage{tabularx}
\usepackage{tabularray}
\usepackage{array}
%\usepackage[french]{babel}
\usepackage[utf8]{inputenc}
\usepackage[T1]{fontenc}
\usepackage[a4paper,left=2cm,right=2cm,top=2.5cm,bottom=2.5cm]{geometry}
\usepackage{graphicx}
\usepackage{colortbl}
\usepackage{hhline}

\usepackage[hidelinks,unicode=false]{hyperref}
\usepackage{setspace}
\usepackage{xcolor}
\usepackage{fontenc}
\usepackage[T1]{fontenc}
\usepackage{titlesec}



% --- Palette technique ---
\definecolor{hdr}{HTML}{1A2E3C}      % bleu nuit --- en-têtes de tableau
\definecolor{rowA}{HTML}{EFF4F7}     % bleu très clair --- ligne paire
\definecolor{rowB}{HTML}{FFFFFF}     % blanc --- ligne impaire
\definecolor{accent}{HTML}{2E6DA4}   % bleu moyen --- titres section
\definecolor{codebg}{HTML}{F4F4F4}  % fond listing
\definecolor{codetxt}{HTML}{1A1A1A} % texte listing

% --- Espacement ---
\setlength{\parindent}{0pt}
\setlength{\parskip}{5pt}
\singlespacing

% --- Titres ---
\titleformat{\section}{\large\bfseries\color{hdr}}{\thesection}{1em}{}[{\color{accent}\titlerule[0.6pt]}]
\titleformat{\subsection}{\normalsize\bfseries\color{hdr}}{\thesubsection}{1em}{}
\titleformat{\subsubsection}{\small\bfseries\color{accent}}{\thesubsubsection}{1em}{}
\titlespacing*{\section}{0pt}{12pt}{6pt}
\titlespacing*{\subsection}{0pt}{8pt}{4pt}
\titlespacing*{\subsubsection}{0pt}{6pt}{3pt}

% --- Macro tableau technique ---
% Colonne gauche : fond sombre, texte blanc gras
% Colonnes données : alternance subtile
\newcommand{\tblhdr}[1]{\cellcolor{hdr}\textcolor{white}{\textbf{#1}}}

\begin{document}
\sloppy

% ════════════════════════════════════════════════════════════════════════════
%  PAGE DE GARDE
% ════════════════════════════════════════════════════════════════════════════
\begin{titlepage}
% [logo université]
\begin{center}
{\large Université Paris Cité --- UFR Mathématiques et Informatique}
\end{center}
\vspace{2.5cm}

\vspace{1.5cm}

\begin{center}
{\Large \textbf{AI Taskforce}}\\[0.4cm]
{\Huge \textbf{Plan de Tests}}\\[0.3cm]
{\small\color{gray} Version 1.0 --- 22 Avril 2026}\\[0.6cm]
{\small Encadrant : M. Janiszek David}
\end{center}
\vspace{2cm}
\begin{center}
\begin{tabular}{ll}
Auteur : & KEITA Fodé Laye \\
N\textsuperscript{o} étudiant : & 22525615 \\
\end{tabular}
\end{center}
\vfill
\begin{center}
\textbf{Membres du projet (Groupe-L3U1)}
\end{center}
\begin{center}
Keita Fodé Laye \quad Guidjou Danil \quad|\quad Dzemchankova Lizaveta \quad|\quad Diallo Aminata \quad|
\end{center}
\end{titlepage}

\renewcommand{\contentsname}{Sommaire}
\tableofcontents
\setcounter{section}{0}
\renewcommand{\thesection}{\arabic{section}}

% ════════════════════════════════════════════════════════════════════════════
\section{Introduction}
% ════════════════════════════════════════════════════════════════════════════

Ce document constitue le plan de tests du projet \textbf{AI Taskforce}.
Il définit la stratégie de validation du système multi-agents et couvre trois niveaux
d'exigence : tests fonctionnels, tests unitaires et tests d'intégration.
Il sert de référence commune entre les développeurs, les testeurs et l'encadrant.


\section{Objectifs et méthodes}
Ce document valide la conformité du produit aux exigences à chaque étape du développement.
Il couvre trois niveaux de tests : unitaires, cohésion et utilisateurs. Chaque scénario
précise le contexte, les acteurs, les prérequis et les résultats attendus.

\section{Documents de référence}
\begin{itemize}
  \item Cahiers de charges et recettes;
  \item modèle de cahier de plan de tests (forge);
  \item comptes rendus de réunions.
\end{itemize}

\section{Guide de lecture}
Selon le profil du lecteur :
\begin{itemize}
  \item \textbf{Développeurs :} vérification de l'implémentation.
  \item \textbf{Testeurs :} validation du fonctionnement global.
  \item \textbf{Enseignant :} évaluation de la conformité au cahier des charges.
\end{itemize}

La \textbf{maîtrise d'œuvre} (groupe projet) assure développement et validation technique.
La \textbf{maîtrise d'ouvrage} (enseignant encadrant) évalue la conformité, qualité et pertinence des solutions.

% ════════════════════════════════════════════════════════════════════════════
\section{Stratégie de test}
% ════════════════════════════════════════════════════════════════════════════

\subsection{Périmètre}

\noindent
\begin{tblr}{
  width = \linewidth,
  colspec = {
    Q[l,3.2cm,bg=hdr,fg=white,font=\bfseries]
    X[l,bg=rowA]
    Q[c,1.8cm,bg=rowA]
    Q[l,3.2cm,bg=rowA]
  },
  hlines = {white},
  vlines = {white},
  rows = {m, rowsep=3pt},
}
Niveau & Périmètre & Tests & Outil \\

Tests unitaires
& Agent, Orchestration, GraphBuilder, requestLLM, RAGService, FASTAPI (agents \& workflows)
& 23
& pytest \\

Tests d'intégration 
& Orchestration$\leftrightarrow$GraphBuilder\newline
  Agent$\leftrightarrow$Pléiade\newline
  FastAPI$\leftrightarrow$PostgreSQL\newline
  Frontend$\leftrightarrow$Backend\newline
  Pipeline RAG 
& 5 
& pytest / TestClient \\

Tests fonctionnels
& Création d'agents, configuration workflow, exécution collaborative, enrichissement RAG, ergonomie, navigation, responsivité
& 14
& Navigateur web \\

\end{tblr}

\subsection{Ordonnancement}

Les niveaux s'exécutent séquentiellement :
\textbf{Unitaires} $\rightarrow$ \textbf{Intégration} $\rightarrow$ \textbf{Fonctionnels}.
Chaque niveau ne démarre que si le niveau précédent est validé.

\subsection{Environnement de test}

\noindent
\begin{tblr}{
  width = \linewidth,
  colspec = {Q[l,4cm,bg=hdr,fg=white,font=\bfseries] X[l,bg=rowA]},
  hlines = {white},
  vlines = {white},
  rows = {m, rowsep=3pt},
}
OS & Windows / Mac OS \\
Python & 3.13 --- environnement virtuel \texttt{.venv} \\
Backend & \texttt{uvicorn} sur \texttt{localhost:8000} \\
Frontend & \texttt{Vite} sur \texttt{localhost:5173} \\
Base de données & PostgreSQL 15 --- instance dédiée aux tests \\
API LLM & Pléiade (\texttt{pleiade.mi.parisdescartes.fr}) \\
Outils de test & pytest, httpx, FastAPI TestClient, unittest.mock \\
\end{tblr}

% ════════════════════════════════════════════════════════════════════════════
\section{Tests fonctionnels}
% ════════════════════════════════════════════════════════════════════════════

Les tests fonctionnels valident les scénarios métier de bout en bout via le navigateur.
Ils couvrent les exigences REQ-AGENT-01, REQ-WKFW-01 et REQ-IHM-01.

% ---
\subsection{TF-AGENT : Création et configuration d'un agent}

\textbf{Exigence :} REQ-AGENT-01 \hfill \textbf{Priorité :} critique

\vspace{4pt}
\noindent
\begin{tblr}{
  width = \linewidth,
  colspec = {Q[l,4cm,bg=hdr,fg=white,font=\bfseries] X[l,bg=rowA]},
  hlines = {white},
  vlines = {white},
  rows = {m, rowsep=3pt},
}
Identifiant & TF-AGENT-01 \\
Description & Créer un agent avec des paramètres valides et vérifier sa persistance. \\
Dépendances & Backend démarré ; base de données accessible. \\
Procédure &
\begin{minipage}[t]{\linewidth}\vspace{2pt}
\begin{enumerate}\setlength\itemsep{1pt}
  \item Accéder à la page de création d'agent.
  \item Renseigner : nom, modèle Pléiade, prompt système, température 0.5, max\_tokens 512.
  \item Soumettre le formulaire et vérifier l'apparition dans la liste.
\end{enumerate}\vspace{2pt}
\end{minipage} \\
Critères d'acceptation & L'agent est persisté (GET \texttt{/agents} le retourne). Les attributs correspondent aux valeurs saisies. \\
Statut & En attente de validation \\
\end{tblr}

\vspace{6pt}
\noindent
\begin{tblr}{
  width = \linewidth,
  colspec = {Q[l,4cm,bg=hdr,fg=white,font=\bfseries] X[l,bg=rowA]},
  hlines = {white},
  vlines = {white},
  rows = {m, rowsep=3pt},
}
Identifiant & TF-AGENT-02 \\
Description & Activer la recherche web (\texttt{use\_web=true}) et vérifier l'enrichissement de la réponse. \\
Dépendances & TF-AGENT-01 \\
Procédure &
\begin{minipage}[t]{\linewidth}\vspace{2pt}
\begin{enumerate}\setlength\itemsep{1pt}
  \item Créer un agent avec l'option \textit{Recherche Web} activée.
  \item Soumettre un prompt requérant une information récente.
  \item Comparer la réponse avec celle d'un agent sans recherche web.
\end{enumerate}\vspace{2pt}
\end{minipage} \\
Critères d'acceptation & La réponse contient des données absentes du modèle de base. Aucune erreur de connexion. \\
Statut & Fonctionnel \\
\end{tblr}

% ---
\subsection{TF-WKFW : Configuration et exécution du workflow}

\textbf{Exigence :} REQ-WKFW-01 \hfill \textbf{Priorité :} critique

\vspace{4pt}
\noindent
\begin{tblr}{
  width = \linewidth,
  colspec = {Q[l,4cm,bg=hdr,fg=white,font=\bfseries] X[l,bg=rowA]},
  hlines = {white},
  vlines = {white},
  rows = {m, rowsep=3pt},
}
Identifiant & TF-WKFW-01 \\
Description & Configurer un workflow multi-agents et vérifier l'exécution collaborative. \\
Dépendances & TF-AGENT-01 ; au moins deux agents créés. \\
Procédure &
\begin{minipage}[t]{\linewidth}\vspace{2pt}
\begin{enumerate}\setlength\itemsep{1pt}
  \item Accéder à l'interface graphique de workflow (ReactFlow).
  \item Glisser-déposer : Planificateur $\rightarrow$ Rédacteur $\rightarrow$ Critique $\rightarrow$ Planificateur.
  \item Valider et enregistrer la configuration.
  \item Soumettre un prompt de test.
\end{enumerate}\vspace{2pt}
\end{minipage} \\
Critères d'acceptation & Workflow persisté en base. Réponse finale produite sans boucle infinie. \\
Statut & Fonctionnel \\
\end{tblr}

\vspace{6pt}
\noindent
\begin{tblr}{
  width = \linewidth,
  colspec = {Q[l,4cm,bg=hdr,fg=white,font=\bfseries] X[l,bg=rowA]},
  hlines = {white},
  vlines = {white},
  rows = {m, rowsep=3pt},
}
Identifiant & TF-WKFW-02 \\
Description & Vérifier le comportement distinct des trois niveaux de recherche. \\
Dépendances & TF-WKFW-01 \\
Procédure &
\begin{minipage}[t]{\linewidth}\vspace{2pt}
\begin{enumerate}\setlength\itemsep{1pt}
  \item Soumettre le même prompt avec \texttt{niveau\_recherche = 1, 2, 3}.
  \item Inspecter les logs : nombre d'appels LLM par sous-tâche.
\end{enumerate}\vspace{2pt}
\end{minipage} \\
Critères d'acceptation &
\begin{minipage}[t]{\linewidth}\vspace{2pt}
\begin{tblr}{width=\linewidth, colspec={Q[c,2cm,bg=rowB,font=\bfseries] X[l,bg=rowB] X[l,bg=rowB]}, hlines={gray!30}, vlines={gray!30}}
Niveau & Appels LLM max & Comportement \\
1 & 1 / sous-tâche & Réponse rapide, qualité de base \\
2 & 2 / sous-tâche & Raffinement superviseur possible \\
3 & 3 / sous-tâche & Meilleure qualité, temps plus long \\
\end{tblr}
\vspace{2pt}
\end{minipage} \\
Statut & Fonctionnel \\
\end{tblr}

\vspace{6pt}
\noindent
\begin{tblr}{
  width = \linewidth,
  colspec = {Q[l,4cm,bg=hdr,fg=white,font=\bfseries] X[l,bg=rowA]},
  hlines = {white},
  vlines = {white},
  rows = {m, rowsep=3pt},
}
Identifiant & TF-WKFW-03 \\
Description & Vérifier un workflow avec enrichissement RAG (document externe). \\
Dépendances & TF-WKFW-01 \\
Procédure &
\begin{minipage}[t]{\linewidth}\vspace{2pt}
\begin{enumerate}\setlength\itemsep{1pt}
  \item Ajouter un document (pdf, docx, txt) à la création d'un agent.
  \item Configurer un workflow contenant l'agent enrichi.
  \item Soumettre une requête nécessitant les données du document.
\end{enumerate}\vspace{2pt}
\end{minipage} \\
Critères d'acceptation & Fichier vectorisé sans erreur. Réponse contient des informations issues du document. \\
Statut & fonctionnel\\
\end{tblr}

\vspace{6pt}
\noindent
\begin{tblr}{
  width = \linewidth,
  colspec = {Q[l,4cm,bg=hdr,fg=white,font=\bfseries] X[l,bg=rowA]},
  hlines = {white},
  vlines = {white},
  rows = {m, rowsep=3pt},
}
Identifiant & TF-WKFW-04 \\
Description & Vérifier la persistance et la recharge d'une configuration de workflow. \\
Dépendances & TF-WKFW-01 \\
Procédure &
\begin{minipage}[t]{\linewidth}\vspace{2pt}
\begin{enumerate}\setlength\itemsep{1pt}
  \item Configurer et enregistrer un workflow.
  \item Recharger la page et ouvrir le workflow sauvegardé.
  \item Vérifier la fidélité du graphe (nœuds, arêtes, paramètres).
\end{enumerate}\vspace{2pt}
\end{minipage} \\
Critères d'acceptation & Configuration restituée intégralement. Workflow éditable après rechargement. \\
Statut & Fonctionnel \\
\end{tblr}

% ---
\subsection{TF-IHM : Ergonomie, navigation et responsivité}

\textbf{Exigence :} REQ-IHM-01 \hfill \textbf{Priorité :} haute

Ces tests valident la qualité de l'interface utilisateur : lisibilité, cohérence de la navigation,
comportement sur différentes résolutions et accessibilité de base.
Ils sont réalisés manuellement via le navigateur (Chrome / Firefox) avec les DevTools ouverts.

\vspace{6pt}
\noindent
\begin{tblr}{
  width = \linewidth,
  colspec = {Q[l,4cm,bg=hdr,fg=white,font=\bfseries] X[l,bg=rowA]},
  hlines = {white},
  vlines = {white},
  rows = {m, rowsep=3pt},
}
Identifiant & TF-IHM-01 \\
Description & Vérifier la navigation entre les pages principales de l'application. \\
Dépendances & Frontend démarré (\texttt{localhost:5173}). \\
Procédure &
\begin{minipage}[t]{\linewidth}\vspace{2pt}
\begin{enumerate}\setlength\itemsep{1pt}
  \item Depuis la page d'accueil, naviguer vers : Dashboard, Création d'agent, Workflow, Exécution.
  \item Utiliser le bouton \textit{Retour} du navigateur sur chaque page.
  \item Vérifier que la barre de navigation est présente et fonctionnelle sur toutes les pages.
\end{enumerate}\vspace{2pt}
\end{minipage} \\
Critères d'acceptation & Toutes les pages sont accessibles sans erreur. Le bouton \textit{Retour} ne provoque pas de page blanche. La navbar est visible et cliquable sur chaque page. \\
Statut & En attente de validation \\
\end{tblr}

\vspace{6pt}
\noindent
\begin{tblr}{
  width = \linewidth,
  colspec = {Q[l,4cm,bg=hdr,fg=white,font=\bfseries] X[l,bg=rowA]},
  hlines = {white},
  vlines = {white},
  rows = {m, rowsep=3pt},
}
Identifiant & TF-IHM-02 \\
Description & Vérifier le chargement et l'affichage du Dashboard sur écran standard (1920$\times$1080). \\
Dépendances & TF-IHM-01 ; utilisateur connecté avec au moins un agent et un workflow. \\
Procédure &
\begin{minipage}[t]{\linewidth}\vspace{2pt}
\begin{enumerate}\setlength\itemsep{1pt}
  \item Ouvrir le Dashboard sur une résolution 1920$\times$1080.
  \item Vérifier l'affichage des 4 tuiles de statistiques, de la liste des agents, des workflows et des documents.
  \item Vérifier que les éléments ne se chevauchent pas et que les textes sont lisibles.
\end{enumerate}\vspace{2pt}
\end{minipage} \\
Critères d'acceptation & Tous les blocs sont visibles et correctement alignés. Aucun débordement horizontal. Textes non tronqués. \\
Statut & En attente de validation \\
\end{tblr}

\vspace{6pt}
\noindent
\begin{tblr}{
  width = \linewidth,
  colspec = {Q[l,4cm,bg=hdr,fg=white,font=\bfseries] X[l,bg=rowA]},
  hlines = {white},
  vlines = {white},
  rows = {m, rowsep=3pt},
}
Identifiant & TF-IHM-03 \\
Description & Vérifier la responsivité de l'interface sur une résolution réduite (768$\times$1024 --- tablette). \\
Dépendances & TF-IHM-01 \\
Procédure &
\begin{minipage}[t]{\linewidth}\vspace{2pt}
\begin{enumerate}\setlength\itemsep{1pt}
  \item Ouvrir les DevTools (F12) et passer en mode responsive à 768$\times$1024.
  \item Naviguer sur les pages : Dashboard, Création d'agent, Workflow.
  \item Vérifier que la mise en page s'adapte (grilles, marges, tailles de texte).
\end{enumerate}\vspace{2pt}
\end{minipage} \\
Critères d'acceptation & Aucun élément ne dépasse la largeur de l'écran. Les boutons restent cliquables. La navbar reste accessible. \\
Statut & En attente de validation \\
\end{tblr}

\vspace{6pt}
\noindent
\begin{tblr}{
  width = \linewidth,
  colspec = {Q[l,4cm,bg=hdr,fg=white,font=\bfseries] X[l,bg=rowA]},
  hlines = {white},
  vlines = {white},
  rows = {m, rowsep=3pt},
}
Identifiant & TF-IHM-04 \\
Description & Vérifier les messages d'erreur et les états vides de l'interface. \\
Dépendances & TF-IHM-01 \\
Procédure &
\begin{minipage}[t]{\linewidth}\vspace{2pt}
\begin{enumerate}\setlength\itemsep{1pt}
  \item Accéder au Dashboard sans aucun agent ni workflow créé.
  \item Soumettre le formulaire de création d'agent avec des champs obligatoires vides.
  \item Simuler une perte de connexion backend (arrêter \texttt{uvicorn}) et recharger le Dashboard.
\end{enumerate}\vspace{2pt}
\end{minipage} \\
Critères d'acceptation & Les états vides affichent un message explicite (ex. : \textit{Aucun agent créé}). Les champs invalides sont signalés visuellement. L'interface ne plante pas en cas d'indisponibilité backend. \\
Statut & En attente de validation \\
\end{tblr}

\vspace{6pt}
\noindent
\begin{tblr}{
  width = \linewidth,
  colspec = {Q[l,4cm,bg=hdr,fg=white,font=\bfseries] X[l,bg=rowA]},
  hlines = {white},
  vlines = {white},
  rows = {m, rowsep=3pt},
}
Identifiant & TF-IHM-05 \\
Description & Vérifier la suppression d'un document depuis le Dashboard. \\
Dépendances & TF-WKFW-03 ; au moins un document indexé visible dans le Dashboard. \\
Procédure &
\begin{minipage}[t]{\linewidth}\vspace{2pt}
\begin{enumerate}\setlength\itemsep{1pt}
  \item Survoler un document dans la section Documents du Dashboard.
  \item Cliquer sur l'icône de suppression et confirmer la boîte de dialogue.
  \item Vérifier la disparition immédiate du document dans la liste.
  \item Vérifier via GET \texttt{/agents/user/\{id\}/documents} que le document n'est plus en base.
\end{enumerate}\vspace{2pt}
\end{minipage} \\
Critères d'acceptation & Le document disparaît de la liste sans rechargement de page. Il est absent en base et ses chunks sont supprimés de ChromaDB. \\
Statut & En attente de validation \\
\end{tblr}

% ════════════════════════════════════════════════════════════════════════════
\section{Tests d'intégration}
% ════════════════════════════════════════════════════════════════════════════

Les tests d'intégration vérifient les interfaces entre modules.
Les dépendances LLM sont mockées sauf mention explicite.

% ---
\subsection{TI-OG-01 : Orchestration $\leftrightarrow$ GraphBuilder}

\textbf{Priorité :} critique \hfill \textbf{Dépendances :} UT-O01, UT-G01

\textbf{Objectif :} vérifier qu'\texttt{orchestration.py} génère une state machine LangGraph
valide à partir d'un JSON de workflow via \texttt{graphBuilder.py}.

\noindent
\begin{tblr}{
  width = \linewidth,
  colspec = {Q[l,4cm,bg=hdr,fg=white,font=\bfseries] X[l,bg=rowA]},
  hlines = {white}, vlines = {white}, rows = {m, rowsep=3pt},
}
Mise en \oe uvre &
  Instancier \texttt{Orchestration} avec \texttt{patern.json}.
  Remplacer le LLM par un mock via \texttt{unittest.mock.patch}.
  Vérifier les n\oe uds du graphe produit par \texttt{graphBuilder}.
  Aucun appel réseau réel. \\
Données d'entrée & \texttt{patern.json} --- LLM mock retournant \texttt{"réponse simulée"}. \\
Résultats attendus & Objet \texttt{Orchestration} créé. N\oe uds et arêtes conformes au JSON. \\
Critères & \texttt{graph.nodes} contient tous les identifiants du JSON. Aucune exception. \\
Statut & En attente de validation \\
\end{tblr}

\subsection{TI-AP-01 : Agent $\leftrightarrow$ API Pléiade}

\textbf{Priorité :} critique \hfill \textbf{Dépendances :} UT-A01

\textbf{Objectif :} vérifier que \texttt{requestLLM.py} envoie une requête conforme au format
OpenAI et reçoit une réponse valide de l'API Pléiade.

\noindent
\begin{tblr}{
  width = \linewidth,
  colspec = {Q[l,4cm,bg=hdr,fg=white,font=\bfseries] X[l,bg=rowA]},
  hlines = {white}, vlines = {white}, rows = {m, rowsep=3pt},
}
Mise en \oe uvre &
  Appeler \texttt{requestLLM.call\_pleiade} directement avec un prompt court
  et un token Pléiade valide issu du fichier \texttt{.env}.
  Inspecter la structure de la réponse reçue (format OpenAI-compatible). \\
Données d'entrée & Prompt : \texttt{"Couleur du ciel en un mot ?"} --- \texttt{max\_tokens=10} --- \texttt{temperature=0.0}. \\
Résultats attendus & Réponse reçue en $<$ 10 s. Structure conforme au format OpenAI. \\
Critères & \texttt{choices[0].message.content} non vide. Code HTTP Pléiade = 200. \\
Statut & En attente de validation \\
\end{tblr}

% ---
\subsection{TI-DB-01 : FastAPI $\leftrightarrow$ PostgreSQL}

\textbf{Priorité :} critique \hfill \textbf{Dépendances :} migrations appliquées

\textbf{Objectif :} vérifier le cycle CRUD complet entre les endpoints FastAPI et la base
PostgreSQL via SQLAlchemy.

\noindent
\begin{tblr}{
  width = \linewidth,
  colspec = {Q[l,4cm,bg=hdr,fg=white,font=\bfseries] X[l,bg=rowA]},
  hlines = {white}, vlines = {white}, rows = {m, rowsep=3pt},
}
Mise en \oe uvre &
  Utiliser le \texttt{TestClient} de FastAPI sur une base PostgreSQL de test
  vidée avant chaque exécution (fixture \texttt{db\_clean}).
  Enchaîner les opérations POST, DELETE, GET et vérifier les codes HTTP
  et l'état de la base après chaque étape. \\
Données d'entrée & POST \texttt{/agents} : \texttt{\{model:"Pleiade", system\_prompt:"Assistant", max\_tokens:500\}}. \\
Résultats attendus & HTTP 201 à la création. HTTP 204 à la suppression. HTTP 404 après suppression. \\
Critères & Identifiant non nul à la création. Agent absent en base après suppression. \\
Statut & En attente de validation \\
\end{tblr}

% ---
\subsection{TI-FB-01 : Frontend React $\leftrightarrow$ Backend via FastAPI}

\textbf{Priorité :} haute \hfill \textbf{Dépendances :} TI-DB-01

\textbf{Objectif :} vérifier l'absence d'erreurs CORS et la bonne désérialisation des
données lors de la communication frontend/backend.

\noindent
\begin{tblr}{
  width = \linewidth,
  colspec = {Q[l,4cm,bg=hdr,fg=white,font=\bfseries] X[l,bg=rowA]},
  hlines = {white}, vlines = {white}, rows = {m, rowsep=3pt},
}
Mise en \oe uvre &
  Test manuel via les DevTools du navigateur (onglet Réseau).
  Créer un workflow dans l'interface React, cliquer sur Enregistrer,
  puis inspecter les requêtes HTTP émises et la console navigateur.
  CORS configuré pour autoriser \texttt{localhost:5173}. \\
Données d'entrée & Workflow créé dans l'interface : 1 superviseur + 2 agents spécialistes. \\
Résultats attendus & POST \texttt{/workflows} retourne HTTP 201. Workflow visible via GET. \\
Critères & Aucune erreur CORS dans la console. Workflow présent en base après enregistrement. \\
Statut & En attente de validation \\
\end{tblr}

% ---
\subsection{TI-RAG-01 : Pipeline RAG complet}

\textbf{Priorité :} haute \hfill \textbf{Dépendances :} UT-A02, TI-DB-01

\textbf{Objectif :} vérifier que \texttt{RAGService} indexe un document et injecte le
contexte pertinent dans le prompt de l'agent et l'evaluation du RAG.

\noindent
\begin{tblr}{
  width = \linewidth,
  colspec = {Q[l,4cm,bg=hdr,fg=white,font=\bfseries] X[l,bg=rowA]},
  hlines = {white}, vlines = {white}, rows = {m, rowsep=3pt},
}
Mise en \oe uvre &
  Instancier \texttt{RagEvaluator} avec le \texttt{RAGService} et un agent
  dont le contexte a été enrichi par un document (PDF, TXT ou CSV).
  Générer un jeu de questions-réponses de référence via \texttt{generer\_dataset},
  puis mesurer la qualité du RAG via les métriques précision, rappel et F1. \\
Données d'entrée & Agent enrichi par un document --- nombre de questions de test paramétrable. \\
Résultats attendus & Valeurs métriques du système RAG : précision, rappel, score F1. \\
Critères & Document correctement indexé. Terme du document présent dans la réponse de l'agent. \\
Statut & Fonctionnel \\
\end{tblr}

% ════════════════════════════════════════════════════════════════════════════
\section{Tests unitaires des modules}
% ════════════════════════════════════════════════════════════════════════════

Chaque module est testé isolément avec \texttt{pytest}.
Toutes les dépendances externes (API Pléiade, base de données) sont mockées
via \texttt{unittest.mock.patch}.

% ---
\subsection{Module \texttt{Agent.py}}

\textbf{Mise en œuvre :} instanciation directe de la classe \texttt{Agent}.
L'API Pléiade est mockée. Les fichiers de test (PDF, CSV) sont dans \texttt{tests/fixtures/}.

\vspace{6pt}
\noindent\textbf{Jeux d'essai}

\vspace{4pt}
\noindent
\begin{tblr}{
  width = \linewidth,
  colspec = {Q[c,2cm,bg=hdr,fg=white,font=\bfseries] X[l,bg=rowA] X[l,bg=rowA]},
  hlines = {white}, vlines = {white}, rows = {m, rowsep=3pt},
}
ID & Données d'entrée & Critères d'évaluation \\
UT-A01 & \texttt{nom="TestAgent"}, \texttt{temperature=0.5}, \texttt{max\_token=512} & \texttt{isinstance(agent, Agent)} vrai. Attributs conformes. \\
UT-A02 & \texttt{tests/fixtures/sample.pdf} (1 page lisible) & \texttt{agent.context} non vide. Pas de \texttt{FileNotFoundError}. \\
UT-A03 & \texttt{tests/fixtures/data.csv} (5 lignes, 3 colonnes) & En-têtes CSV présents dans le contexte. \\
\end{tblr}

% ---
\subsection{Module \texttt{orchestration.py}}

\textbf{Mise en œuvre :} instanciation d'\texttt{Orchestration} avec superviseur, spécialistes
et niveau de recherche. Le LLM est remplacé par un mock à réponse prédéterminée.
On valide \texttt{executer()}.

\vspace{6pt}
\noindent
\begin{tblr}{
  width = \linewidth,
  colspec = {Q[c,2cm,bg=hdr,fg=white,font=\bfseries] X[l,bg=rowA] X[l,bg=rowA]},
  hlines = {white}, vlines = {white}, rows = {m, rowsep=3pt},
}
ID & Données d'entrée & Critères d'évaluation \\
UT-O01 & \texttt{niveau=1}, prompt court, LLM mock & 1 appel LLM/sous-tâche. Réponse de type \texttt{str}. \\
UT-O02 & \texttt{niveau=2}, prompt complexe, LLM mock & $\leq$ 2 appels LLM/sous-tâche. \\
UT-O03 & \texttt{niveau=3}, prompt multi-facettes, LLM mock & $\leq$ 3 appels LLM/sous-tâche. \\
UT-O04 & Prompt : \texttt{"\ \ \ "} (espaces uniquement) & \texttt{pytest.raises(ValueError, match="prompt vide")}. \\
\end{tblr}

% ---
\subsection{Module \texttt{graphBuilder.py}}

\textbf{Mise en œuvre :} passage de JSON valides et invalides à \texttt{graphBuilder}.
Inspection des attributs du graphe LangGraph produit.

\vspace{6pt}
\noindent
\begin{tblr}{
  width = \linewidth,
  colspec = {Q[c,2cm,bg=hdr,fg=white,font=\bfseries] X[l,bg=rowA] X[l,bg=rowA]},
  hlines = {white}, vlines = {white}, rows = {m, rowsep=3pt},
}
ID & Données d'entrée & Critères d'évaluation \\
UT-G01 & \texttt{patern.json} (fichier de workflow valide) & \texttt{graph.nodes} contient tous les IDs du JSON. Aucune exception. \\
UT-G02 & \texttt{\{"nodes": []\}} (pas de superviseur) & \texttt{ValueError} ou \texttt{KeyError} levée. \\
\end{tblr}

% ---
\subsection{Module \texttt{requestLLM.py}}

\textbf{Mise en œuvre :} interception des appels réseau via \texttt{responses} ou
\texttt{httpretty} pour simuler les réponses de l'API Pléiade.

\vspace{6pt}
\noindent
\begin{tblr}{
  width = \linewidth,
  colspec = {Q[c,2cm,bg=hdr,fg=white,font=\bfseries] X[l,bg=rowA] X[l,bg=rowA]},
  hlines = {white}, vlines = {white}, rows = {m, rowsep=3pt},
}
ID & Données d'entrée & Critères d'évaluation \\
UT-R01 & Prompt court, \texttt{max\_tokens=50}, mock HTTP 200 format OpenAI & \texttt{choices[0].message.content} non vide. \\
UT-R02 & Mock Pléiade retournant HTTP 401 & \texttt{AuthenticationError} ou équivalent levée. \\
\end{tblr}

% ---
\subsection{Module \texttt{RAGService.py}}

\textbf{Mise en œuvre :} instanciation directe de \texttt{RAGService}.
Les fichiers de test (PDF, TXT, CSV) sont dans \texttt{tests/fixtures/}.
Le modèle d'embeddings est utilisé localement ; aucun appel à l'API Pléiade.

\vspace{6pt}
\noindent\textbf{Jeux d'essai}

\vspace{4pt}
\noindent
\begin{tblr}{
  width = \linewidth,
  colspec = {Q[c,2cm,bg=hdr,fg=white,font=\bfseries] X[l,bg=rowA] X[l,bg=rowA]},
  hlines = {white}, vlines = {white}, rows = {m, rowsep=3pt},
}
ID & Données d'entrée & Critères d'évaluation \\
UT-RAG01 & \texttt{tests/fixtures/sample.pdf} (document lisible, contenu connu) & Document indexé sans erreur. Nombre de chunks $>$ 0. \\
UT-RAG02 & Requête sémantique correspondant à un passage connu du document & Chunk pertinent retourné en premier résultat. Similarité $>$ seuil configuré. \\
UT-RAG03 & Fichier vide (0 octet) & \texttt{ValueError} ou \texttt{EmptyDocumentError} levée. Aucun chunk en base. \\
UT-RAG04 & Requête sans correspondance dans le document indexé & Liste de résultats vide ou score sous le seuil. Aucune exception. \\
UT-RAG05 & Document déjà indexé soumis une seconde fois & Pas de doublon en base. Nombre de chunks identique après réindexation. \\
\end{tblr}

% ---
\subsection{API REST --- \texttt{agent\_router.py}}

\textbf{Mise en œuvre :} \texttt{TestClient} FastAPI avec base PostgreSQL de test réinitialisée
avant chaque suite.

\vspace{6pt}
\noindent
\begin{tblr}{
  width = \linewidth,
  colspec = {Q[c,2.2cm,bg=hdr,fg=white,font=\bfseries] X[l,bg=rowA] X[l,bg=rowA]},
  hlines = {white}, vlines = {white}, rows = {m, rowsep=3pt},
}
ID & Données d'entrée & Critères d'évaluation \\
UT-API01 & POST \texttt{/agents} --- \texttt{\{model, system\_prompt, max\_tokens\}} valides & HTTP 201. Champ \texttt{id} non nul. \\
UT-API02 & GET \texttt{/agents/\{id\}} avec un identifiant existant & HTTP 200. Corps JSON conforme au schéma agent. \\
UT-API03 & GET \texttt{/agents/\{id\}} avec un identifiant inexistant & HTTP 404. Message d'erreur explicite. \\
UT-API04 & DELETE \texttt{/agents/\{id\}} avec un identifiant existant & HTTP 204. Agent absent en base après suppression. \\
UT-API05 & POST \texttt{/agents} avec corps invalide (\texttt{max\_tokens} négatif) & HTTP 422. Détail de validation présent dans la réponse. \\
\end{tblr}

% ---
\subsection{API REST --- \texttt{workflow\_router.py}}

\textbf{Mise en œuvre :} \texttt{TestClient} FastAPI avec base PostgreSQL de test réinitialisée
avant chaque suite. Les agents nécessaires sont créés en fixture.

\vspace{6pt}
\noindent
\begin{tblr}{
  width = \linewidth,
  colspec = {Q[c,2.2cm,bg=hdr,fg=white,font=\bfseries] X[l,bg=rowA] X[l,bg=rowA]},
  hlines = {white}, vlines = {white}, rows = {m, rowsep=3pt},
}
ID & Données d'entrée & Critères d'évaluation \\
UT-WF01 & POST \texttt{/workflows} --- JSON valide (superviseur + 2 agents, arêtes correctes) & HTTP 201. Champ \texttt{id} non nul. Structure persistée en base. \\
UT-WF02 & GET \texttt{/workflows/\{id\}} avec un identifiant existant & HTTP 200. Nœuds et arêtes conformes au JSON soumis. \\
UT-WF03 & POST \texttt{/workflows} --- JSON sans nœud superviseur & HTTP 422 ou \texttt{ValueError}. Message d'erreur explicite. \\
UT-WF04 & POST \texttt{/workflows} --- arêtes référençant des agents inexistants & HTTP 404. Workflow non persisté en base. \\
\end{tblr}

% ════════════════════════════════════════════════════════════════════════════
\section{Références}
% ════════════════════════════════════════════════════════════════════════════

\begin{description}

    \item[Cahier des charges :] \textit{« AI Taskforce »} --- Document de spécification du projet, Groupe-L3U1 --- Université Paris Cité, 2026.

    \item[Modèle plan de tests :] \textit{« Modèle de plan de tests »} --- David Janiszek, Université Paris Descartes, ref. D7, v1.01.

    \item[Cahier de recettes :] \textit{« Cahier de recettes v1.2.0 »} --- KEITA Fodé Laye, 22 février 2026, Université Paris Cité.

    \item[Documentation FastAPI :] \textit{« FastAPI --- Modern, fast web framework for building APIs »} --- Documentation officielle du framework backend utilisé. \\ \url{https://fastapi.tiangolo.com}

    \item[Documentation LangGraph :] \textit{« LangGraph --- Build stateful, multi-actor applications with LLMs »} --- Documentation officielle du framework d'orchestration multi-agents. \\ \url{https://langchain-ai.github.io/langgraph/}

    \item[Documentation pytest :] \textit{« pytest --- helps you write better programs »} --- Documentation officielle du framework de tests unitaires et d'intégration. \\ \url{https://docs.pytest.org}

    \item[Documentation ReactFlow :] \textit{« React Flow --- Wire Your Ideas with React Flow »} --- Documentation officielle de la bibliothèque de visualisation de graphes utilisée pour l'interface workflow. \\ \url{https://reactflow.dev}

    \item[Documentation ChromaDB :] \textit{« Chroma --- the AI-native open-source embedding database »} --- Documentation officielle de la base vectorielle utilisée pour le RAG. \\ \url{https://docs.trychroma.com}

    \item[Documentation LangChain :] \textit{« LangChain --- Build context-aware reasoning applications »} --- Documentation officielle des composants LangChain utilisés (SemanticChunker, retrievers). \\ \url{https://python.langchain.com/docs}

\end{description}

\end{document}
